\chapter*{Power Series and Taylor Series}
\addcontentsline{toc}{chapter}{Power Series and Taylor Series}

\section*{Power Series}

A power series centered at $a$ has the form
$$\sum_{n=0}^{\infty} c_n(x-a)^n.$$ 
For each fixed $x$, this becomes a numerical series, so convergence depends on $x$.

\begin{definition}[Radius and interval of convergence]
There exists a number $R\in[0,\infty]$ such that the power series converges for
$$|x-a|<R$$
and diverges for
$$|x-a|>R.$$ 
The number $R$ is the \textbf{radius of convergence}. Endpoints must be tested separately.
\end{definition}

In many problems, the ratio test is the easiest way to find $R$.

\section*{Differentiation and Integration of Power Series}

Inside the interval of convergence, power series behave like polynomials term by term:
\begin{align*}
\frac{d}{dx}\sum_{n=0}^{\infty} c_n(x-a)^n &= \sum_{n=1}^{\infty} n c_n(x-a)^{n-1}, \\
\int \sum_{n=0}^{\infty} c_n(x-a)^n\,dx &= C + \sum_{n=0}^{\infty} \frac{c_n}{n+1}(x-a)^{n+1}.
\end{align*}

This makes power series an efficient tool for producing new expansions from known ones.

\section*{Taylor and Maclaurin Series}

Suppose $f$ can be represented near $x=a$ by a power series
$$f(x)=\sum_{n=0}^{\infty} c_n(x-a)^n.$$ 
Differentiating repeatedly and evaluating at $x=a$ gives
$$c_n = \frac{f^{(n)}(a)}{n!}.$$ 
So we obtain the Taylor series:

\begin{definition}[Taylor series]
The Taylor series of $f$ centered at $a$ is
$$f(x) = \sum_{n=0}^{\infty} \frac{f^{(n)}(a)}{n!}(x-a)^n,$$
provided the series converges to $f(x)$.
\end{definition}

When $a=0$, this is called the \textbf{Maclaurin series}.

\section*{Important Standard Expansions}

The most frequently used series are:
\begin{align*}
\frac{1}{1-x} &= \sum_{n=0}^{\infty} x^n, \qquad |x|<1, \\
 e^x &= \sum_{n=0}^{\infty} \frac{x^n}{n!}, \\
\sin x &= \sum_{n=0}^{\infty} (-1)^n\frac{x^{2n+1}}{(2n+1)!}, \\
\cos x &= \sum_{n=0}^{\infty} (-1)^n\frac{x^{2n}}{(2n)!}, \\
\ln(1+x) &= \sum_{n=1}^{\infty} (-1)^{n-1}\frac{x^n}{n}, \qquad -1<x\le 1.
\end{align*}

\section*{Taylor Polynomials}

The degree-$n$ Taylor polynomial at $a$ is
$$T_n(x)=\sum_{k=0}^{n}\frac{f^{(k)}(a)}{k!}(x-a)^k.$$ 
It gives a local polynomial approximation to $f$ near $a$.

\begin{remark}
In practice, Taylor polynomials are useful because complicated functions can be approximated by low-degree polynomials, which are much easier to compute and analyze.
\end{remark}